\documentclass[10pt,a4paper]{amsart}
\usepackage[margin=19mm]{geometry}
\usepackage{amsmath,amssymb,mathtools}
\usepackage[scr=rsfs,cal=euler]{mathalfa}
\usepackage{xfrac}
\usepackage[unicode,hidelinks,bookmarks=false]{hyperref}
\newcommand{\ie}{i.e.\ }
\newcommand{\cf}{cf.\ }
\newcommand{\resp}{respectively\ }
\newcommand{\Subsection}{\subsection}
\providecommand{\nobreakdash}{\nobreak\hbox{-}}
\setlength{\emergencystretch}{2em}
\multlinegap0pt
\usepackage{amssymb}
\usepackage{enumerate}
\usepackage{algorithm}\usepackage{algpseudocode}
\usepackage{tikz}
\usepackage{mathtools}
\newtheorem{assumption}{Assumption}
\newtheorem{lemma}{Lemma}
\newcommand{\Fm}{F_m}
\title{Stochastic block coordinate and function alternation for multi-objective optimization and learning}
\author{Trang H. Tran, Luis Nunes Vicente}
\date{Source: arXiv:2605.12432v1 (CC BY 4.0)}
\begin{document}
\maketitle

\noindent\emph{Source and licence.} Stochastic block coordinate and function alternation for multi-objective optimization and learning. Version arXiv:2605.12432v1 (CC BY 4.0), distributed by arXiv under the Creative Commons Attribution 4.0 International licence, http://creativecommons.org/licenses/by/4.0/, as recorded in the arXiv licence metadata for this exact version and shown on the abstract page https://arxiv.org/abs/2605.12432v1. Full licence notice: \texttt{licenses/arxiv-2605.12432-CC-BY-4.0.txt}.\par\smallskip

\noindent\textit{Editorial scope.} Counted unit: the article's lemma lem:bounds (source0.tex lines 753--760). The article's complete original proof of it is delivered with it (source0.tex lines 761--792, one \texttt{proof} environment of 32 lines). No mathematics is written, repaired or re-derived: the statement and the proof are the article's own lines, in the article's own notation, and no retained line is substituted or re-flowed.  Retained uncounted as the source-local context the proof needs: lines 384--393; lines 749--751.  Nothing was omitted from any retained span, so the package carries no omission record.  No retained line contains a citation, so the wrapper carries no bibliography and no bibliography entry.  No retained line contains a figure environment, an image-inclusion command or drawing code, so no image is delivered and the package declares no asset dependency.  Editorial formatting only, in addition: an AMS \texttt{amsart} preamble that restates the article's own declarations and macros used by the retained lines, namely the environments \texttt{assumption}, \texttt{lemma} on separate counters, together with the standard packages those lines need.  The retained environments print in this document as Assumption 1, Lemma 1 in order of appearance, whereas the article prints the same items with the numbering of its own full document.  The complete native source of the article as distributed in the arXiv e-print for this exact version is bundled unchanged as \texttt{sources/200422/source0.tex}.
\begin{assumption}[Unbiasedness and Bounded Gradient]\label{assump:conds}
For every $k \in \{ 1, \dots, q\}$, every $x \in \mathbb{R}^n$,
and every realization of $\xi$, the stochastic gradient $\nabla g_k(x,\xi)$ satisfies the following conditions
\begin{itemize}
\item[(a)] %Unbiasedness: 
$\mathbb{E}_{\xi} \left[\nabla g_k(x,\xi)\right] = \nabla f_k(x)$.
\item[(b)] %Bounded gradient: 
$\mathbb{E}_{\xi}[\|\nabla g_k(x,\xi)\|^2] \leq \sigma^2$, where $\sigma^2$ is a constant. 
\end{itemize} 
\end{assumption}

\begin{align}\label{eq:alg_update}
    x^{t,i,j+1} = x^{t,i,j} - \alpha_t U_{\sigma(i+1)} \nabla_{\sigma(i+1)} g_{\pi(j)} (x^{t,i,j}, \xi^{t,i,j}).
\end{align}

\begin{lemma}\label{lem:bounds}
Let Assumption \ref{assump:conds} holds. 
We have 
\begin{itemize}
\item[(a)] For every $x\in \mathbb{R}^n$, $\|\nabla f_j(x)\| \leq \sigma$ and $\|\nabla \Fm(x)\| \leq \sigma$.
\item[(b)] For every $t, i,$ and $j$, $\mathbb{E} \left\|x^{t,i,j+1}-x^{t,i,j}\right\| \leq \alpha_t \sigma $ and $\mathbb{E} \left[\left\|x^{t,i,j+1}-x^{t,i,i}\right\|^2\right] \leq \alpha^2_t \sigma^2$.
\end{itemize} 
\end{lemma}

\begin{proof} 
\begin{itemize}
\item[(a)] We obtain the first bound by applying the unbiased property in Assumption \ref{assump:conds}(a), the Jensen's inequality, and the bounded gradient condition in Assumption \ref{assump:conds}(b),
\begin{align}\label{eq:bound_grad}
    \|\nabla f_j(x)\| =  \|\mathbb{E}_{\xi} \left[\nabla g_j(x,\xi)\right]\| \leq \sqrt{ \mathbb{E}_{\xi} \left[\|\nabla g_j(x,\xi)\|^2\right]} \leq \sigma.
\end{align}
The second bound follows from the definition of $\Fm$ and the triangle inequality, i.e.,
\begin{align*}
    \|\nabla \Fm(x)\| =  \left\|\frac{1}{p}\sum_{j=0}^{p-1}\nabla f_{\pi(j)} (x)\right\| \leq \frac{1}{p}\sum_{j=0}^{p-1}\|\nabla f_{\pi(j)} (x)\| \overset{\eqref{eq:bound_grad}}{\leq} \sigma.
\end{align*}
%By convexity of the norm operator, $\|\nabla \Fm(x)\|  \leq \sigma$.
\item[(b)] 
From equation~\eqref{eq:alg_update}, we have 
\begin{align*}
    \|x^{t,i,j+1} - x^{t,i,j}\| = \alpha_t \| U_{\sigma(i+1)} \nabla_{\sigma(i+1)} g_{\pi(j)} (x^{t,i,j}, \xi^{t,i,j})\|.
\end{align*}
Applying expectation with respect to $\mathcal{F}_{t,i,j}$, then using Jensen's inequality, the fact that $\| U_{\sigma(i+1)}\| \leq 1$, and Assumption \ref{assump:conds}(b), one obtains 
\begin{align*}
    \mathbb{E} [\|x^{t,i,j+1} - x^{t,i,j}\|  | \mathcal{F}_{t,i,j}]&= \alpha_t  \mathbb{E} [\| U_{\sigma(i+1)} \nabla_{\sigma(i+1)} g_{\pi(j)} (x^{t,i,j}, \xi^{t,i,j})\|| \mathcal{F}_{t,i,j}]\\
    &\leq \alpha_t  \sqrt{\mathbb{E} [\| U_{\sigma(i+1)} \nabla_{\sigma(i+1)} g_{\pi(j)} (x^{t,i,j}, \xi^{t,i,j})\|^2| \mathcal{F}_{t,i,j}]} \leq \alpha_t \sigma.
\end{align*}
Similar arguments yield
\begin{align*}
    \mathbb{E} [\|x^{t,i,j+1} - x^{t,i,j}\|^2  | \mathcal{F}_{t,i,j}]&= \alpha_t^2  \mathbb{E} [\| U_{\sigma(i+1)} \nabla_{\sigma(i+1)} g_{\pi(j)} (x^{t,i,j}, \xi^{t,i,j})\|^2| \mathcal{F}_{t,i,j}]\leq \alpha_t^2 \sigma^2.
\end{align*}
Taking total expectation, we obtain
\begin{align*}
    \mathbb{E} \left\|x^{t,i,j+1}-x^{t,i,j}\right\| \leq \alpha_t \sigma \text{ and }
    \mathbb{E} \left[\left\|x^{t,i,j+1}-x^{t,i,j}\right\|^2\right] \leq \alpha_t^2 \sigma^2.
\end{align*}
\end{itemize} 
\end{proof}

\end{document}
